\documentclass[11pt]{article}
\usepackage[margin=1in]{geometry}
\usepackage[T1]{fontenc}
\usepackage[utf8]{inputenc}
\usepackage{lmodern}
\usepackage{amsmath}
\usepackage{booktabs}
\usepackage{hyperref}

\title{From Candidate Patches to an Automatable Refactoring Engine}
\author{rmatch optimization workflow note}
\date{\today}

\begin{document}
\maketitle

\begin{abstract}
This note documents the optimization loop currently used in the \texttt{rmatch} repository: generate a candidate patch, evaluate feasibility, enforce correctness gates, and accept only candidates that show measurable improvement in a fixed A/B performance gate (10K patterns, 1MB corpus). The same loop can be generalized into a ``refactoring engine'' that proposes, validates, and ranks code transformations with reproducible evidence.
\end{abstract}

\section{Current Loop (As Practiced)}
The current process is intentionally conservative and evidence-driven:
\begin{enumerate}
  \item \textbf{Generate candidate patch.} Form a concrete implementation hypothesis (for example, reducing allocations or hash lookups in a hot path).
  \item \textbf{Feasibility screening.} Perform static inspection to reject ideas that are too invasive, semantically risky, or mismatched to known bottlenecks.
  \item \textbf{Correctness gate loop.} Run fast unit and smoke checks, fix regressions, and repeat until the candidate is stable.
  \item \textbf{Performance gate.} Run a fixed local A/B benchmark: baseline branch vs candidate branch on the same 10K/1MB workload.
  \item \textbf{Decision rule.} Keep the candidate only when measured runtime improves (or at least does not regress) with acceptable variance and no correctness regressions.
\end{enumerate}

\section{Why This Works}
The loop is effective because it separates claims:
\begin{itemize}
  \item \textbf{Correctness claim:} proven by tests and smoke checks.
  \item \textbf{Performance claim:} proven by controlled A/B measurement.
  \item \textbf{Merge decision:} based on evidence, not intuition.
\end{itemize}
This reduces ``optimization theater'' where plausible code changes fail to improve real workloads.

\section{Formalization}
Let \(c_i\) be candidate \(i\). Let \(G_c(c_i)\in\{0,1\}\) be correctness-gate pass/fail and \(R(c_i)\) be runtime ratio (candidate over baseline) from the 10K/1MB A/B gate.

A practical accept policy is:
\[
\text{Accept}(c_i) \iff G_c(c_i)=1 \;\wedge\; R(c_i)\le \tau
\]
where \(\tau\) is a configurable threshold (for example \(\tau=1.00\) for strict non-regression, or slightly above 1.00 when noise bounds are explicitly modeled).

\section{Generalizing to a Refactoring Engine}
The same structure can be automated end-to-end:
\begin{enumerate}
  \item \textbf{Candidate generator.} Produces patch proposals from patterns (allocation reduction, data-structure substitution, branch simplification, caching).
  \item \textbf{Static risk estimator.} Scores semantic risk and expected effort.
  \item \textbf{Execution orchestrator.} Runs the correctness loop and benchmark gates in isolated branches/environments.
  \item \textbf{Evidence store.} Persists tests, benchmark artifacts, ratios, and decision outcomes.
  \item \textbf{Policy engine.} Applies acceptance thresholds and prioritizes follow-up experiments.
\end{enumerate}

At scale, this becomes a high-performance refactoring system: a closed loop where each proposal is validated against objective correctness and efficiency criteria before promotion.

\section{Near-Term Extensions}
To evolve from the current manual loop into a robust engine:
\begin{itemize}
  \item Add richer gates (memory, p95/p99 latency, larger workload cohorts).
  \item Add automatic rollback and retry strategies for flaky measurements.
  \item Calibrate noise-aware thresholds per hardware cohort.
  \item Build a historical model linking patch features to observed gains.
\end{itemize}

\section{Conclusion}
The current \texttt{rmatch} optimization workflow is already a practical self-improvement loop. Documenting it as a formal pipeline makes automation straightforward and creates a foundation for a future refactoring engine that can continuously search for safe, measurable performance improvements.

\end{document}
